\documentclass{article}
\usepackage[pdftex]{graphicx}
\begin{document}

\author{The clm-mem project\\ \small (working repository: \texttt{clm-mem}, v0.3.0)}
\date{\today}

\begin{document}
\maketitle
\begin{abstract}
\noindent
Context Language Models (CLMs) treat the context as an editable file that the model itself manages, and report state-of-the-art accuracy at lower compute than external context-management baselines. However, the headline CLM implementation rewrites the whole context every request, which invalidates the KV cache on each turn: a cost paid either by novel serving machinery or by reduced throughput.
We present \textbf{clm-mem}, a zero-training memory harness that keeps the two ideas we consider essential from CLMs---(i)~the model classifies and writes its own context, and (ii)~context is a mirrored on-disk artifact---while replacing free-form context editing with a small deterministic eviction engine.
The same chat model that reads and writes text also tags its own context as \emph{long-term}, \emph{short-term} or \emph{transient} while decoding; the harness then enforces a token budget in three monotone phases (drop transient, FIFO rounds on short-term, offload long-term to an archive).
Because compaction happens only when the window reaches its maximum size and never rewrites mid-turn, KV cache churn is reduced without requiring new cache strategies; any improvement to KV management helps both approaches.
We describe the software architecture implemented in \texttt{clm\_mem} (blocks, inline and post-turn classification, a typed decision backstop, three-phase eviction, an offload-don't-destroy archive, session-scoped serving with workspaces and an embedded MCP toolbox), define a marking protocol that turns harness events into labels for scoring tag quality, and design the benchmarking program: synthetic stress regimes already run against the deterministic stub backend (budget never violated; $100\%$ of emitted facts remain recoverable from the archive), plus planned evaluations on ContextBench-style tasks, LoCoMo, LongMemEval and RULER with ablations over the classification mode, the FIFO knob $\kappa$, and the backstop.
\end{abstract}
\newpage
\end{document}
